\documentclass[10pt,a4paper]{amsart}
\usepackage[margin=19mm]{geometry}
\usepackage{amsmath,amssymb,mathtools}
\usepackage[scr=rsfs,cal=euler]{mathalfa}
\usepackage{xfrac}
\usepackage[unicode,hidelinks,bookmarks=false]{hyperref}
\newcommand{\ie}{i.e.\ }
\newcommand{\cf}{cf.\ }
\newcommand{\resp}{respectively\ }
\newcommand{\Subsection}{\subsection}
\providecommand{\nobreakdash}{\nobreak\hbox{-}}
\setlength{\emergencystretch}{2em}
\multlinegap0pt
\usepackage{amssymb}
\usepackage{tikz}
\newtheorem{lemma}{Lemma}
\title{On lower bounds for the number of ideal and finite vertices of right-angled hyperbolic polyhedra in dimensions from 5 to 12}
\author{Andrey Egorov}
\date{Source: arXiv:2603.29925v1 (CC BY 4.0)}
\begin{document}
\maketitle

\noindent\emph{Source and licence.} On lower bounds for the number of ideal and finite vertices of right-angled hyperbolic polyhedra in dimensions from 5 to 12. Version arXiv:2603.29925v1 (CC BY 4.0), distributed by arXiv under the Creative Commons Attribution 4.0 International licence, http://creativecommons.org/licenses/by/4.0/, as recorded in the arXiv licence metadata for this exact version and shown on the abstract page https://arxiv.org/abs/2603.29925v1. Full licence notice: \texttt{licenses/arxiv-2603.29925-CC-BY-4.0.txt}.\par\smallskip

\noindent\textit{Editorial scope.} Counted unit: the article's lemma lem:fin\_recurrence (source0.tex lines 301--305). The article's complete original proof of it is delivered with it (source0.tex lines 306--308, one \texttt{proof} environment of 3 lines). No mathematics is written, repaired or re-derived: the statement and the proof are the article's own lines, in the article's own notation, and no retained line is substituted or re-flowed.  No further source-local context is needed: every cross-reference of every retained line resolves inside the two delivered spans.  Nothing was omitted from any retained span, so the package carries no omission record.  No retained line contains a citation, so the wrapper carries no bibliography and no bibliography entry.  No retained line contains a figure environment, an image-inclusion command or drawing code, so no image is delivered and the package declares no asset dependency.  Editorial formatting only, in addition: an AMS \texttt{amsart} preamble that restates the article's own declarations and macros used by the retained lines, namely the environments \texttt{lemma} on separate counters, together with the standard packages those lines need.  The retained environments print in this document as Lemma 1 in order of appearance, whereas the article prints the same items with the numbering of its own full document.  The complete native source of the article as distributed in the arXiv e-print for this exact version is bundled unchanged as \texttt{sources/197555/source0.tex}.
	\begin{lemma} \label{lem:fin_recurrence}
		Let $P^n \in \mathcal{P}^n$. Then the minimal number of its finite vertices is bounded from below by the relation:
		$$v_{fin}(P^n) \ge \left\lceil \frac{a_{n-1}(P^n) \cdot v'_{fin}(P^n)}{n} \right\rceil$$
		where $v'_{fin}(P^n)$ is the minimal number of finite vertices in a facet of the polyhedron $P^n$.
	\end{lemma}

	\begin{proof}
		Summing the number of finite vertices over all facets, we obtain $v_{fin}(P^n) \cdot n = \sum_{\dim F = n-1} v_{fin}(F) \ge a_{n-1}(P^n) \cdot v'_{fin}(P^n)$. Expressing $v_{fin}(P^n)$ and taking into account the integrality of the number of vertices, we arrive at the required inequality.
	\end{proof}

\end{document}
